\begin{proposition}{4.2}\label{XIII.4.2}
Let $r$ be the nilpotent rank of $\mathfrak g$, and let
$a\in\mathfrak g$. Then one has
\[
 \operatorname{rank}_k\operatorname{Nil}(a,\mathfrak g)\geq r,
\]
and one has equality if and only if one has
\[
 c_{n-r}(a)\ne0.
\]
In this case, $\operatorname{Nil}(a,\mathfrak g)$ is a nilpotent Lie
algebra (and we shall see in 5.7 b) the converse, when $\mathfrak g$
is the Lie algebra of a smooth algebraic group $G$ over $k$).
\end{proposition}
The first assertion is trivial, for by definition one has
$\operatorname{rank}_k\operatorname{Nil}(a,\mathfrak g)=$ the multiplicity
of the zero root in $P_{\mathfrak g}(a,t)$. Let us prove that if
$c_{n-r}(a)\ne0$, then $\operatorname{Nil}(a)$ is nilpotent, which also
means that for every $x\in\operatorname{Nil}(a)$,
$\ad_{\operatorname{Nil}(a)}(x)$ is a nilpotent endomorphism. One may
suppose $k$ algebraically closed; then, since
$\ad(a)_{\mathfrak g/\operatorname{Nil}(a)}$ is injective, there exists
a nonempty open subset $U$ of $\mathrm W(\operatorname{Nil}(a))$ such
that for every $x\in U(k)$, $\ad(x)_{\mathfrak g/\operatorname{Nil}(a)}$
is injective, hence $\operatorname{Nil}(x)\subset\operatorname{Nil}(a)$;
one may moreover suppose $U$ contained in the open subset of points
where $c_{n-r}$ does not vanish (since this open subset is nonempty by
$c_{n-r}(a)\ne0$), and then $\operatorname{Nil}(x)$ having the same
dimension as $\operatorname{Nil}(a)$, one will have
$\operatorname{Nil}(x)=\operatorname{Nil}(a)$. Consequently, for every
$x\in U(k)$, $\ad(x)_{\operatorname{Nil}(a)}$ is nilpotent, and by the
principle of extension of algebraic identities, this will remain true
for every $x\in\operatorname{Nil}(a)$, so $\operatorname{Nil}(a)$ is nilpotent.

One says that the element $a$ of $\mathfrak g$ is \emph{regular} if
$c_{n-r}(a)\ne0$, i.e. if $\operatorname{rank}_k\operatorname{Nil}(a,\mathfrak g)=r$.
When $k$ is infinite, this therefore also means that
$\operatorname{rank}_k\operatorname{Nil}(a,\mathfrak g)$ is as small
as possible (for $a$ varying in $\mathfrak g$). In all cases, the notion
of regular element of $\mathfrak g$ is invariant under extension of
the base field, and the set of points of $\mathrm W(\mathfrak g)$ which
are regular (i.e. which come from regular points of
$\mathrm W(\mathfrak g)$ with values in a suitable extension of $k$)
is open, for it is identical to $\mathrm W(\mathfrak g)_{c_{n-r}}$ (the
set of points where $c_{n-r}$ is invertible).
% typo: unmatched opening parenthesis after rank_k removed.

\begin{corollary}{4.3}\label{XIII.4.3}
Let $a$ be a regular element of $\mathfrak g$, and $\mathfrak h$ a Lie
subalgebra of $\mathfrak g$ containing $a$. Then $\mathfrak h$ is
nilpotent if and only if $\mathfrak h\subset\operatorname{Nil}(a,\mathfrak g)$;
in particular, $\operatorname{Nil}(a,\mathfrak g)$ is a maximal
nilpotent subalgebra of $\mathfrak g$.
\end{corollary}
Since $\operatorname{Nil}(a)$ is nilpotent, the relation
$\mathfrak h\subset\operatorname{Nil}(a)$ indeed implies that
$\mathfrak h$ is nilpotent, and conversely, if $\mathfrak h$ is
nilpotent, it is contained in the nilspace of its element $a$, i.e.
$\mathfrak h\subset\operatorname{Nil}(a)$.

\begin{proposition}{4.4}\label{XIII.4.4}
Suppose $k$ infinite. Let $\mathfrak d$ be a Lie subalgebra of
$\mathfrak g$. Consider the following conditions:
\begin{enumerate}[label=\textup{(\roman*)}]
\item $\mathfrak d$ is maximal nilpotent and contains a regular element
of $\mathfrak g$.
\item[\textup{(i bis)}] $\mathfrak d$ is of the form
$\operatorname{Nil}(a,\mathfrak g)$, where $a$ is a regular element of
$\mathfrak g$.
\item $\mathfrak d$ is nilpotent and of the form
$\operatorname{Nil}(a,\mathfrak g)$, where $a\in\mathfrak g$.
\item[\textup{(ii bis)}] $\mathfrak d$ is nilpotent, and there exists
$a\in\mathfrak d$ such that $\ad(a)_{\mathfrak g/\mathfrak d}$ is injective.
\item $\mathfrak d$ is nilpotent and identical to its own normalizer.
\end{enumerate}
One has the implications:
\[
 \textup{(i)}\Leftrightarrow\textup{(i bis)}\Rightarrow\textup{(ii)}
 \Leftrightarrow\textup{(ii bis)}\Leftrightarrow\textup{(iii)}
\]
(and we shall see in 5.7 a) that if $\mathfrak g$ is the Lie algebra
of a smooth algebraic group, then all the preceding conditions are
equivalent).
\end{proposition}
The equivalence of (i) and (i bis) is trivial by 4.3, and these
conditions trivially imply (ii). The equivalence of (ii) and (ii bis)
is also trivial, as is (ii bis) $\Rightarrow$ (iii) (cf. 4.1). It remains
to prove implication (iii) $\Rightarrow$ (ii bis), the only one,
moreover, which uses the fact that $k$ is infinite, and which follows
at once from the

\begin{lemma}{4.5}\label{XIII.4.5}
Let $\mathfrak d$ be a nilpotent Lie algebra over an infinite field
$k$, acting on a finite-dimensional vector space $V$. Suppose that
for every $x\in\mathfrak d$, the endomorphism $u(x)$ is not injective.
Then there exists a nonzero element $v$ of $V$ annihilated by
$\mathfrak d$.
\end{lemma}
One may suppose $k$ algebraically closed and $\mathfrak d$
finite-dimensional. One then knows that $V$ is a direct sum of finitely
many nonzero stable vector subspaces $V_i$ ($1\leq i\leq n$), such
that for every $i$, and every $x\in\mathfrak d$, $u(x)|_{V_i}$ has a
single eigenvalue $\lambda_i(x)$ (cf. Bourbaki, \emph{Groupes et
Alg\`ebres de Lie}, Chap. I, \S 4, Exercise 22). Let $c_i(x)$ be the
constant term of the characteristic polynomial of $u(x)|_{V_i}$, so
that $\lambda_i(x)=0$ if and only if $c_i(x)=0$. Then $c_i$ is a
polynomial function on $\mathfrak d$, and the hypothesis means that
$\mathfrak d$ is the union of the zero sets of the $c_i$. Thus one
of the $c_i$ is zero, which reduces one (replacing $V$ by $V_i$) to
the case where $V$ is such that the $u(x)$ ($x\in\mathfrak d$) are
nilpotent. But then Engel's theorem (Bourbaki \loccit th. 1) implies
that there exists a nonzero $v$ in $V$ annihilated by $\mathfrak d$.
\hfill Q.E.D.

One easily sees that ($k$ still being an infinite field) conditions
(i), (i bis) of 4.4 are invariant under every extension of the base
field. If they hold, one will say that $\mathfrak d$ is a \emph{Cartan
subalgebra} of $\mathfrak g$; in the general case ($k$ not necessarily
infinite) one will say that $\mathfrak d$ is a Cartan subalgebra of
$\mathfrak g$ if it becomes a Cartan subalgebra for one (and consequently
every) extension of the base field $k\to k'$, with $k'$ infinite.
This therefore implies that $\mathfrak d$ is nilpotent and equal
to its own normalizer.

\begin{proposition}{4.6}\label{XIII.4.6}
\noindent a) Let $a$ be an element of $\mathfrak g$. If $a$ is regular,
it is contained in one and only one Cartan subalgebra of $\mathfrak g$
(and we shall see in 6.1 d) the converse when $k$ is algebraically
closed and $\mathfrak g$ the Lie algebra of a smooth algebraic group).

\noindent b) Let $\mathfrak d$ be a Cartan subalgebra of $\mathfrak g$,
and $a$ an element of $\mathfrak d$; then $a$ is regular in
$\mathfrak g$ if and only if $\ad(a)_{\mathfrak g/\mathfrak d}$ is injective.
\end{proposition}
Indeed, for a) one notes that if $a$ is regular, then
$\operatorname{Nil}(a,\mathfrak g)$ is a Cartan subalgebra of
$\mathfrak g$ (for this is true over an infinite extension $k'$ of
$k$), and it then follows at once from 4.3 that every Cartan subalgebra
of $\mathfrak g$ containing $a$ is identical to the preceding one.
For b) one notes that the nullity\nde{That is, the dimension of the
nilspace.} of $\ad(a)_{\mathfrak g}$ is equal to the sum of the
nullities of $\ad(a)_{\mathfrak d}$ and of
$\ad(a)_{\mathfrak g/\mathfrak d}$, and since the first is $r$, the
sum is equal to $r$ if and only if $\ad(a)_{\mathfrak g/\mathfrak d}$
is injective. \hfill Q.E.D.

\begin{corollary}{4.7}\label{XIII.4.7}
Let $a$ be a regular element of $\mathfrak g$, $\mathfrak d$ a Cartan
subalgebra of $\mathfrak g$ containing $a$, $A$ an algebra over $k$,
$\mathfrak g_A$ and $\mathfrak d_A\subset\mathfrak g_A$ the $A$-Lie
algebras obtained from $\mathfrak g,\mathfrak d$ by base change,
and $a_A$ the image of $a$ in $\mathfrak g_A$. Let $u$ be an automorphism
of $\mathfrak g_A$. For $u(\mathfrak d_A)=\mathfrak d_A$ to hold,
it is necessary and sufficient that one have $u(a_A)\in\mathfrak d_A$.
\end{corollary}
The condition is trivially necessary; let us prove that it is also
sufficient. If it holds, then $\mathfrak d'=u(\mathfrak d_A)$ is a
Lie subalgebra containing $a_A$, and every element $b$ of it is such
that $\ad(b)_{\mathfrak d}$ is nilpotent (for $\mathfrak d'$ is
isomorphic to $\mathfrak d_A$, which has this property, as follows
at once from the definition of ``nilpotent'' in Bourbaki, \emph{Groupes
et Alg\`ebres de Lie}, Chap. I, \S 4, def. 1). Taking $b=a_A$, one
sees that the nilspace $\operatorname{Nil}(b,\mathfrak g_A)$ contains
$\mathfrak d'$; on the other hand, it is equal to $\mathfrak d_A$,
and since $\mathfrak d'$ is locally a direct summand in the module
$\mathfrak g_A$ ($\mathfrak d$ being so), hence in $\mathfrak d_A$,
and it is a projective module of the same rank $r$ as the latter,
one concludes that it is equal to it. \hfill Q.E.D.
% typo?: kept as printed: d'=u(d_A) contains a_A, and ad(b)_d.

\begin{proposition}{4.8}\label{XIII.4.8}
Let $\mathfrak h$ be a Lie subalgebra of $\mathfrak g$.

\noindent a) The following conditions are equivalent if $k$ is infinite:
\begin{enumeratei}
\item $\mathfrak h$ contains a Cartan subalgebra $\mathfrak d$ of
$\mathfrak g$.
\item $\mathfrak h$ contains a regular element $a$ of $\mathfrak g$,
and an element $b$ such that $\ad(b)_{\mathfrak g/\mathfrak h}$ is injective.
\item $\mathfrak h$ has the same nilpotent rank as $\mathfrak g$,
and contains a regular element of $\mathfrak g$.
\end{enumeratei}
These conditions are invariant under extension of the base field $k$.

\noindent b) Suppose that these conditions hold over a suitable infinite
extension $k'$ of $k$. Let $a\in\mathfrak h$; then $a$ is regular in
$\mathfrak g$ if and only if it is regular in $\mathfrak h$ and
$\ad(a)_{\mathfrak g/\mathfrak h}$ is injective, i.e. if and only if
$\operatorname{Nil}(a,\mathfrak h)=\operatorname{Nil}(a,\mathfrak g)$
and this is a Cartan subalgebra of $\mathfrak h$.

\noindent c) Under the conditions of b), let $\mathfrak d$ be a Lie
subalgebra of $\mathfrak h$; for it to be a Cartan subalgebra of
$\mathfrak h$, it suffices that it be a Cartan subalgebra of
$\mathfrak g$ (and we shall see in 5.8 that the condition is also
necessary if $\mathfrak g$ is the Lie algebra of a smooth algebraic
group $G$, and $\mathfrak h$ the Lie algebra of a smooth algebraic
subgroup $H$ of $G$).
\end{proposition}
One sees immediately that conditions (ii) and (iii) of a) are
invariant under extension of the base field $k$ (supposed infinite),
and that in statements b) and c), one may suppose $k$ infinite,
which we shall do. If $\mathfrak h$ contains the Cartan subalgebra
$\mathfrak d=\operatorname{Nil}(a,\mathfrak g)$, then $a$ is a regular
element of $\mathfrak g$ such that $\ad(a)_{\mathfrak g/\mathfrak h}$
is injective, hence (i) $\Rightarrow$ (ii). Conversely, if (ii) holds,
then for a ``sufficiently general'' element $a$ of $\mathfrak h$,
$a$ simultaneously satisfies the two conditions considered in (ii),
hence $\operatorname{Nil}(a,\mathfrak g)$ is a Cartan subalgebra of
$\mathfrak g$ and is contained in $\mathfrak h$, so one has (i).
Thus (i) and (ii) are equivalent. Suppose they hold, and let $a$ be
a variable element of $\mathfrak h$; then
\[
\begin{aligned}
 \operatorname{rank}_k\operatorname{Nil}(a,\mathfrak g)
 &=\operatorname{rank}_k\operatorname{Nil}(a,\mathfrak h)\\
 &\quad+\operatorname{rank}_k\operatorname{Nil}(\ad(a)_{\mathfrak g/\mathfrak h}),
\end{aligned}\tag{x}
\]
on the other hand, the two terms on the right-hand side are respectively
$\geq r'=$ the nilpotent rank of $\mathfrak h$, and $\geq0$, the
equalities moreover being attained\nde{``The equality'' has been
replaced by ``the equalities''.} for a ``sufficiently general''
element of $\mathfrak h$. Moreover, one also has
$\operatorname{rank}_k\operatorname{Nil}(a,\mathfrak g)\geq r=$ the
nilpotent rank of $\mathfrak g$, equality being attained for a
``sufficiently general'' element of $\mathfrak h$, and being attained
if and only if $a$ is regular in $\mathfrak g$. One concludes from
this that one has $r=r'$, and that $a$ is regular if and only if
the two terms on the right-hand side of (x) are equal respectively
to $r'$ and to $0$, i.e. if and only if $a$ is regular in
$\mathfrak h$ and $\ad(a)_{\mathfrak g/\mathfrak h}$ is injective,
which proves b); and c) follows trivially by taking an element $a$
in $\mathfrak d$ regular in $\mathfrak g$, so that
$\operatorname{Nil}(a,\mathfrak g)=\mathfrak d$. Moreover, the preceding
result shows that (i) $\Rightarrow$ (iii); finally (iii)
$\Rightarrow$ (i), for under (iii), a sufficiently general element
$a$ of $\mathfrak h$ is regular in $\mathfrak h$ and in
$\mathfrak g$, so $\operatorname{Nil}(a,\mathfrak h)\subset
\operatorname{Nil}(a,\mathfrak g)$ are respectively Cartan subalgebras
of $\mathfrak h$ and of $\mathfrak g$, and since they have the same
rank over $k$, they are identical, which proves (i). This completes
the proof of 4.8.

\sgasection{5}{Case of the Lie algebra of a smooth algebraic group: density theorem}
Let $G$ be a smooth algebraic group over the field $k$, and
$\mathfrak g$ its Lie algebra. Let $\mathfrak h$ be a Lie subalgebra
of $\mathfrak g$. Let $G$ act on $\mathrm W(\mathfrak g)$ by the
adjoint representation, and consider the subscheme
$\mathrm W(\mathfrak h)$. The construction of \No 1 leads us to introduce
\[
 N=\Norm_G(\mathfrak h)=\Norm_G(\mathrm W(\mathfrak h)),
\]
which is an algebraic subgroup of $G$ (not necessarily smooth),
\[
 \mathfrak n=\Lie(N),
\]
and the scheme
\[
 X=G\times^N\mathrm W(\mathfrak h)
\]
the quotient of $X=G\times\mathrm W(\mathfrak h)$ by $N$ acting on
the right by $(g,x)\cdot n=(gn,\Ad(n^{-1})x)$.
% typo?: kept as printed: X is used again for G times W(h).
We consider the canonical morphisms
\[
\xymatrix{
 G\times\mathrm W(\mathfrak h)\ar[d]_q\ar[dr]^\varphi&\\
 X\ar[r]_\psi&\mathrm W(\mathfrak g).
}
\]

\begin{theorem}{5.1}\label{XIII.5.1}
With the preceding notation, suppose $k$ infinite. Consider the
following conditions:
\begin{enumeratei}
\item $\mathfrak h$ contains a Cartan subalgebra $\mathfrak d$ of
$\mathfrak g$.
\item There exists $a\in\mathfrak h$ such that
$\ad(a)_{\mathfrak g/\mathfrak h}$ is injective.
\item $\varphi:G\times\mathrm W(\mathfrak h)\to\mathrm W(\mathfrak g)$
is generically smooth.
\item The preceding morphism $\varphi$ (or again
$\psi:X\to\mathrm W(\mathfrak g)$) is dominant, and $\mathfrak h$
has the same nilpotent rank as $\mathfrak g$.
\item $\psi:X\to\mathrm W(\mathfrak g)$ is generically smooth and
$\mathfrak h=\mathfrak n$.
\item $\psi:X\to\mathrm W(\mathfrak g)$ is dominant, and
$\mathfrak h=\mathfrak n$.
\item $\psi:X\to\mathrm W(\mathfrak g)$ is dominant.
\item $\psi:X\to\mathrm W(\mathfrak g)$ is dominant, and $N$ is smooth.
\item $\psi:X\to\mathrm W(\mathfrak g)$ is dominant, and
$\mathfrak h=\mathfrak n$, and $\mathfrak h$ is the Lie algebra of
a smooth algebraic subgroup $H$ of $G$.
\item $\psi:X\to\mathrm W(\mathfrak g)$ is generically quasi-finite,
and $\mathfrak n=\mathfrak h$.
\item $\psi:X\to\mathrm W(\mathfrak g)$ is generically \'etale,
and $\mathfrak n=\mathfrak h$.
\item There exists a smooth algebraic subgroup $H$ of $G$ with
Lie algebra $\mathfrak h$, and $\mathfrak h$ contains a Cartan
subalgebra $\mathfrak d$ of $\mathfrak g$.
\item There exists $a\in\mathfrak h$ such that $N$ and the transporter
$M_a$ of $a$ into $\mathfrak h$ (cf. \No 1) coincide in a neighborhood
of $e$, and $\mathfrak n=\mathfrak h$.
\end{enumeratei}
One then has the following diagram of implications:
\[
\xymatrix@C=1em@R=2.8em{
 (\mathrm{xi})\ar@{<=>}[rr]&&
 (\mathrm{xii})\ar@{<=>}[r]\ar@{=>}[d]&
 (\mathrm{xiii})\ar@{=>}[r]&
 (\mathrm{viii})\ar@{<=>}[r]&
 (\mathrm{ix})\ar@{<=>}[r]\ar@{=>}[d]&(\mathrm{x})\\
 (\mathrm{i})\ar@{<=>}[r]&
 (\mathrm{ii})\ar@{<=>}[r]&
 (\mathrm{iii})\ar@{<=>}[r]&
 (\mathrm{iv})\ar@{=>}[r]&
 (\mathrm{v})\ar@{=>}[r]&
 (\mathrm{vi})\ar@{=>}[r]&(\mathrm{vii}).
}
\]
When $k$ has characteristic zero, all the conditions considered are
equivalent. Finally, one has
\[
 (\mathrm{xi})\Leftrightarrow[(\mathrm{i})\text{ and }(\mathrm{viii})]
 \Leftrightarrow[(\mathrm{v})\text{ and }(\mathrm{viii})].
\]
\end{theorem}
\begin{proof}
Note first the trivial implications:
\[
 (\mathrm{v})\Rightarrow(\mathrm{vi})\Rightarrow(\mathrm{vii}),\qquad
 (\mathrm{ix})\Rightarrow(\mathrm{vi}),\qquad
 (\mathrm{xi})\Leftrightarrow[(\mathrm{x})\text{ and }(\mathrm{v})].
\]
Let us prove the equivalence of conditions (i) to (iv) and the fact
that they imply (v). Implication (i) $\Rightarrow$ (ii) is trivial.
On the other hand, (iii) means, when $k$ is algebraically closed,
that there exists a point of $G\times\mathrm W(\mathfrak h)$ rational
over $k$ at which the tangent map to $\varphi$ is surjective, and one
sees at once that one can take this point of the form $(e,a)$,
where $a\in\mathfrak h$ (transforming it, if necessary, by an
operation of $G(k)$). One concludes that if $k$ is infinite (not
necessarily algebraically closed) this condition (clearly sufficient)
of generic smoothness is still necessary. Now the tangent map is
easily calculated: identifying the tangent space to
$\mathrm W(\mathfrak h)$ at $a$ with $\mathfrak h$, it is the map
\[
 (\xi,x)\mto[\xi,a]+x
\]
from $\mathfrak g\times\mathfrak h$ to $\mathfrak g$. Saying that it
is surjective also means that $\ad(a)_{\mathfrak g/\mathfrak h}$ is
surjective, or, what amounts to the same thing, injective. This
proves the equivalence of conditions (ii) and (iii). Moreover,
(ii) clearly implies $\mathfrak h=\mathfrak n$, and (iii) implies
that $\psi$ is generically smooth, for if $\varphi$ is smooth at a
point $u$, it follows ($q$ being flat) that $\psi$ is smooth at
$q(u)$. Thus (ii), (iii) imply (v). Let us prove that they imply
(i). For this, note that since $\psi$ is dominant, and the set of
regular points of $\mathfrak g$ is dense open, it follows that
$\mathfrak h$ contains regular elements of $\mathfrak g$, so that
a ``sufficiently general'' element $b$ of $\mathfrak h$ is regular
in $\mathfrak g$ and satisfies $\ad(b)_{\mathfrak g/\mathfrak h}$
injective, hence $\operatorname{Nil}(b,\mathfrak g)\subset\mathfrak h$,
so $\mathfrak h$ contains the Cartan subalgebra
$\mathfrak d=\operatorname{Nil}(b,\mathfrak g)$. Thus (i), (ii), (iii)
are equivalent; finally (i) $\Leftrightarrow$ (iv), for one has
already noted that if $\mathfrak h$ contains a Cartan subalgebra,
it has the same rank as $\mathfrak g$ (4.6), hence (i)
$\Rightarrow$ (iv); conversely, if (iv) holds, then $\mathfrak h$
contains a regular element of $\mathfrak g$, and since it has the
same rank as $\mathfrak g$, it contains a Cartan subalgebra by 4.6.

Let us prove the equivalence of conditions (viii) to (x). Note
first the following facts:
\begin{lemma}{5.2}\label{XIII.5.2}
\noindent a) If $N$ is smooth, then
\[
 \dim X\leq\dim G,
\]
equality holding if and only if $\mathfrak h=\mathfrak n$.

\noindent b) If $\mathfrak h=\mathfrak n$, then
\[
 \dim X\geq\dim G,
\]
equality holding if and only if $N$ is smooth.
\end{lemma}
These assertions follow at once from the formula
\[
 \dim X=\dim G-\dim N+\dim_k\mathfrak h,
\]
and from the fact that $\dim N=\dim_k\mathfrak n$ is equivalent to
$N$ smooth.

This being established, (viii) $\Rightarrow$ (x), for (viii) implies
$\dim X\geq\dim G$, hence, by 5.2 a), equality of these dimensions
and $\mathfrak h=\mathfrak n$, hence (x), and one similarly sees
(x) $\Rightarrow$ (viii) by applying 5.2 b). On the other hand,
(viii) implies (ix), for it implies $\mathfrak h=\mathfrak n$,
so $\mathfrak h$ is the Lie algebra of the smooth algebraic subgroup
$N$ of $G$, and conversely (ix) $\Rightarrow$ (viii), for since $H$
normalizes its Lie algebra $\mathfrak h$, and it is contained in $N$,
and since $H$ is smooth and has the same Lie algebra as $N$, it
follows that $N$ is smooth.

Finally, let us prove the equivalence of conditions (xi), (xii),
(xiii) and the fact that they imply (iii) (which will finish
establishing our diagram of implications). One has (xi)
$\Leftrightarrow$ (xiii), for if $\mathfrak n=\mathfrak h$, then
by 5.2 b) one has $\dim X\geq\dim G$, hence (xi) is then equivalent
(taking into account that $\mathrm W(\mathfrak g)$ is normal) to
the fact that $\psi$ is generically unramified, which is also
equivalent to (xiii) by (1.1 (ii) $\Leftrightarrow$ (iii)), proceeding
as above for the proof of (ii) $\Leftrightarrow$ (iii). Since
(xi) $\Rightarrow$ (x) $\Rightarrow$ (viii) by what we have already
seen, one sees that (xi) implies that $N$ is smooth, i.e.
$q:G\times\mathrm W(\mathfrak h)\to X$ is smooth, so the composite
$\varphi=\psi\circ q$ is generically smooth, i.e. one has (iii).
Since (iii) $\Rightarrow$ (i), it also follows that (xi)
$\Rightarrow$ (xii)\nde{For $H=N$ will do.}. Finally (xii)
$\Rightarrow$ (xi), for one clearly has (xii) $\Rightarrow$ (i),
so, since one has seen (i) $\Rightarrow$ (iii) $\Rightarrow$ (v),
one has (xii) $\Rightarrow$ (v); it follows that one also has
(xii) $\Rightarrow$ (ix), and since one has seen (ix)
$\Rightarrow$ (x), it follows that (xii) $\Rightarrow$ ((v) and (x)),
hence (xii) $\Rightarrow$ (xi), since generically \'etale $=$
generically smooth $+$ generically quasi-finite.

Finally, when $k$ has characteristic $0$, then (vii) $\Rightarrow$
(viii), for by a theorem of Cartier, $N$ is automatically smooth
(VI\textsubscript{B} 1.6.1.), and [(viii) and (x)] $\Rightarrow$ (xi),
since in characteristic zero, for a morphism of integral preschemes,
generically \'etale $=$ dominant and generically quasi-finite.
This shows that in this case, all conditions (i) to (xiii) are
equivalent.
\end{proof}

% original p. 288
\begin{corollary}{5.3}\label{XIII.5.3}
Under the equivalent conditions (viii) to (x), there exists a unique
smooth connected algebraic subgroup $H$ of $G$ whose Lie algebra is
$\mathfrak h$, and one has
\[
 \Norm_G(H)=\Norm_G(\mathfrak h)=N,\qquad H=N^0.
\]
\end{corollary}
% The proof of Corollary 5.3 starts in en-05.tex.
